% Options for packages loaded elsewhere
\PassOptionsToPackage{unicode}{hyperref}
\PassOptionsToPackage{hyphens}{url}
\PassOptionsToPackage{table}{xcolor}
%
\documentclass[
  letterpaper,
]{article}
\usepackage{amsmath,amssymb}
\usepackage{iftex}
\usepackage[T1]{fontenc}
\usepackage[utf8]{inputenc}
\usepackage{textcomp}
\usepackage{lmodern}
\ifPDFTeX\else
  % xetex/luatex font selection
\fi
% Use upquote if available, for straight quotes in verbatim environments
\IfFileExists{upquote.sty}{\usepackage{upquote}}{}
\IfFileExists{microtype.sty}{% use microtype if available
  \usepackage[]{microtype}
  \UseMicrotypeSet[protrusion]{basicmath} % disable protrusion for tt fonts
}{}
\makeatletter
\@ifundefined{KOMAClassName}{% if non-KOMA class
  \IfFileExists{parskip.sty}{%
    \usepackage{parskip}
  }{% else
    \setlength{\parindent}{0pt}
    \setlength{\parskip}{6pt plus 2pt minus 1pt}}
}{% if KOMA class
  \KOMAoptions{parskip=half}}
\makeatother
\usepackage{xcolor}
\usepackage{color}
\usepackage{fancyvrb}
\newcommand{\VerbBar}{|}
\newcommand{\VERB}{\Verb[commandchars=\\\{\}]}
\DefineVerbatimEnvironment{Highlighting}{Verbatim}{commandchars=\\\{\}}
% Add ',fontsize=\small' for more characters per line
\newenvironment{Shaded}{}{}
\newcommand{\AlertTok}[1]{\textcolor[rgb]{1.00,0.00,0.00}{\textbf{#1}}}
\newcommand{\AnnotationTok}[1]{\textcolor[rgb]{0.38,0.63,0.69}{\textbf{\textit{#1}}}}
\newcommand{\AttributeTok}[1]{\textcolor[rgb]{0.49,0.56,0.16}{#1}}
\newcommand{\BaseNTok}[1]{\textcolor[rgb]{0.25,0.63,0.44}{#1}}
\newcommand{\BuiltInTok}[1]{\textcolor[rgb]{0.00,0.50,0.00}{#1}}
\newcommand{\CharTok}[1]{\textcolor[rgb]{0.25,0.44,0.63}{#1}}
\newcommand{\CommentTok}[1]{\textcolor[rgb]{0.38,0.63,0.69}{\textit{#1}}}
\newcommand{\CommentVarTok}[1]{\textcolor[rgb]{0.38,0.63,0.69}{\textbf{\textit{#1}}}}
\newcommand{\ConstantTok}[1]{\textcolor[rgb]{0.53,0.00,0.00}{#1}}
\newcommand{\ControlFlowTok}[1]{\textcolor[rgb]{0.00,0.44,0.13}{\textbf{#1}}}
\newcommand{\DataTypeTok}[1]{\textcolor[rgb]{0.56,0.13,0.00}{#1}}
\newcommand{\DecValTok}[1]{\textcolor[rgb]{0.25,0.63,0.44}{#1}}
\newcommand{\DocumentationTok}[1]{\textcolor[rgb]{0.73,0.13,0.13}{\textit{#1}}}
\newcommand{\ErrorTok}[1]{\textcolor[rgb]{1.00,0.00,0.00}{\textbf{#1}}}
\newcommand{\ExtensionTok}[1]{#1}
\newcommand{\FloatTok}[1]{\textcolor[rgb]{0.25,0.63,0.44}{#1}}
\newcommand{\FunctionTok}[1]{\textcolor[rgb]{0.02,0.16,0.49}{#1}}
\newcommand{\ImportTok}[1]{\textcolor[rgb]{0.00,0.50,0.00}{\textbf{#1}}}
\newcommand{\InformationTok}[1]{\textcolor[rgb]{0.38,0.63,0.69}{\textbf{\textit{#1}}}}
\newcommand{\KeywordTok}[1]{\textcolor[rgb]{0.00,0.44,0.13}{\textbf{#1}}}
\newcommand{\NormalTok}[1]{#1}
\newcommand{\OperatorTok}[1]{\textcolor[rgb]{0.40,0.40,0.40}{#1}}
\newcommand{\OtherTok}[1]{\textcolor[rgb]{0.00,0.44,0.13}{#1}}
\newcommand{\PreprocessorTok}[1]{\textcolor[rgb]{0.74,0.48,0.00}{#1}}
\newcommand{\RegionMarkerTok}[1]{#1}
\newcommand{\SpecialCharTok}[1]{\textcolor[rgb]{0.25,0.44,0.63}{#1}}
\newcommand{\SpecialStringTok}[1]{\textcolor[rgb]{0.73,0.40,0.53}{#1}}
\newcommand{\StringTok}[1]{\textcolor[rgb]{0.25,0.44,0.63}{#1}}
\newcommand{\VariableTok}[1]{\textcolor[rgb]{0.10,0.09,0.49}{#1}}
\newcommand{\VerbatimStringTok}[1]{\textcolor[rgb]{0.25,0.44,0.63}{#1}}
\newcommand{\WarningTok}[1]{\textcolor[rgb]{0.38,0.63,0.69}{\textbf{\textit{#1}}}}
\usepackage{longtable,booktabs,array}
\usepackage{calc} % for calculating minipage widths
% Correct order of tables after \paragraph or \subparagraph
\usepackage{etoolbox}
\makeatletter
\patchcmd\longtable{\par}{\if@noskipsec\mbox{}\fi\par}{}{}
\makeatother
% Allow footnotes in longtable head/foot
\IfFileExists{footnotehyper.sty}{\usepackage{footnotehyper}}{\usepackage{footnote}}
\makesavenoteenv{longtable}
\setlength{\emergencystretch}{3em} % prevent overfull lines
\providecommand{\tightlist}{%
  \setlength{\itemsep}{0pt}\setlength{\parskip}{0pt}}
\setcounter{secnumdepth}{5}
% Self-contained presentation rules embedded in every generated LaTeX file.
\usepackage{iftex}
\ifPDFTeX\else\PackageError{SSE-ASAS}{Compile with pdfLaTeX only}{}\fi
\usepackage[T1]{fontenc}
\usepackage[utf8]{inputenc}
\usepackage{lmodern}
\usepackage{microtype}
\usepackage{geometry}
\geometry{margin=0.93in,headheight=16pt}
\usepackage{xcolor} % [table] passed at top of file (option-clash fix)
\definecolor{ink}{HTML}{15283B}
\definecolor{brand}{HTML}{075B75}
\definecolor{brandlight}{HTML}{EAF5F8}
\definecolor{warm}{HTML}{A15C1D}
\definecolor{muted}{HTML}{536577}
\usepackage{fancyhdr}
\usepackage{titlesec}
\usepackage{enumitem}
\usepackage{booktabs}
\usepackage{longtable}
\usepackage{array}
\usepackage{tabularx}
\usepackage{tcolorbox}
\tcbuselibrary{breakable,skins}
\usepackage{listings}
\usepackage{fancyvrb}
\usepackage{hyperref}
\hypersetup{colorlinks=true,linkcolor=brand,urlcolor=brand,citecolor=brand}
\urlstyle{same}
\usepackage{lastpage}
\setlength{\emergencystretch}{3em}
\setlength{\parskip}{0.45em}
\setlength{\parindent}{0pt}
\setlist[itemize]{leftmargin=*,topsep=3pt,itemsep=2pt}
\setlist[enumerate]{leftmargin=*,topsep=3pt,itemsep=2pt}
\titleformat{\section}{\Large\bfseries\color{ink}}{\thesection}{0.75em}{}[\vspace{-2pt}\color{brand}\titlerule]
\titleformat{\subsection}{\large\bfseries\color{brand}}{\thesubsection}{0.7em}{}
\titleformat{\subsubsection}{\normalsize\bfseries\color{ink}}{\thesubsubsection}{0.7em}{}
\pagestyle{fancy}
\fancyhf{}
\fancyhead[L]{\small\sffamily\textcolor{brand}{SOFTWARE SECURITY ENGINEERING}}
\fancyhead[R]{\small\sffamily\textcolor{muted}{TECHNICAL SERIES}}
\fancyfoot[L]{\scriptsize\sffamily Version 1.0 \textbullet\ 14 October 2026}
\fancyfoot[R]{\scriptsize\sffamily Page \thepage\ of \pageref{LastPage}}
\renewcommand{\headrulewidth}{0.5pt}
\renewcommand{\footrulewidth}{0.2pt}
\setcounter{tocdepth}{2}
\setcounter{secnumdepth}{3}
\lstdefinestyle{sse}{basicstyle=\ttfamily\footnotesize,frame=single,framerule=0.4pt,rulecolor=\color{brand!40},backgroundcolor=\color{brandlight!45},numbers=left,numberstyle=\scriptsize\color{muted},numbersep=6pt,breaklines=true,breakatwhitespace=false,tabsize=2,showstringspaces=false,columns=fullflexible,keepspaces=true,captionpos=b,aboveskip=9pt,belowskip=9pt}
\lstset{style=sse}
\newtcolorbox{sourceconcept}[1][]{enhanced,breakable,colback=brandlight,colframe=brand!65!black,title={ORIGINAL CONCEPT},fonttitle=\bfseries\sffamily,#1}
\newtcolorbox{modernextension}[1][]{enhanced,breakable,colback=white,colframe=warm!85!black,title={MODERN EXTENSION},fonttitle=\bfseries\sffamily,#1}
\newtcolorbox{fieldwork}[1][]{enhanced,breakable,colback=white,colframe=ink!55,title={PRACTICAL APPLICATION},fonttitle=\bfseries\sffamily,#1}
\AtBeginDocument{\color{ink}}
\makeatletter
\renewcommand{\maketitle}{%
 \begin{titlepage}\thispagestyle{empty}
 \vspace*{1.5cm}{\color{brand}\rule{\textwidth}{3pt}}\par\vspace{1cm}
 {\sffamily\bfseries\color{muted}\large SOFTWARE SECURITY ENGINEERING\\ APPLICATION SECURITY ASSESSMENT SERIES}\par\vspace{1.05cm}
 {\sffamily\bfseries\color{ink}\fontsize{27}{34}\selectfont\@title\par}\vspace{0.8cm}
 {\sffamily\large\color{brand}\@author\par}\vfill
 \begin{tcolorbox}[colback=brandlight,colframe=brand!55,boxrule=0.5pt]
 \small\textbf{Editorial policy}\quad Original teaching synthesis with book-grounded concepts distinguished from contemporary practices. Modern platform behavior is version-sensitive; validate against linked vendor documentation before production use.
 \end{tcolorbox}
 {\color{brand}\rule{\textwidth}{2pt}}\par\vspace{0.3cm}
 {\sffamily\small Version 1.0 \quad|\quad October 2026 \quad|\quad pdfLaTeX}\end{titlepage}\setcounter{page}{1}}
\makeatother
\ifLuaTeX
  \usepackage{selnolig}  % disable illegal ligatures
\fi
\usepackage{bookmark}
\IfFileExists{xurl.sty}{\usepackage{xurl}}{} % add URL line breaks if available
\urlstyle{same}
\hypersetup{
  pdftitle={Secure Coding and Defensive Implementation},
  pdfauthor={Independent Technical Study Edition},
  hidelinks,
  pdfcreator={LaTeX via pandoc}}

\title{Secure Coding and Defensive Implementation}
\author{Independent Technical Study Edition}
\date{}

\begin{document}
\maketitle

{
\setcounter{tocdepth}{3}
\tableofcontents
}
\section{Purpose and quality
standard}\label{purpose-and-quality-standard}

Secure coding means preserving explicit security invariants across input
validation, authorization, state transitions, memory management,
dependencies, error paths, and deployment boundaries. A patch is
complete when its relevant negative tests pass and the same flaw cannot
be reached through an equivalent entry point.

\textbf{Original basis:} Dowd et al., Chapters 5--13 and 17--18; Chess
and West, Chapters 5--12; Bishop, Chapters 14--18. \textbf{Modern
extensions:} language ecosystems, contemporary framework APIs, OWASP
ASVS 5.0.0, CWE, CodeQL and supply-chain controls. Code is illustrative
and must be adapted to the selected framework and API contract.

\section{Security invariants in application
code}\label{security-invariants-in-application-code}

\subsection{Establish the trusted
context}\label{establish-the-trusted-context}

External data includes HTTP parameters, messages, files, environment
variables, process output, database content originating with users, and
data from third-party APIs. Validate schema, type, length, range,
encoding, provenance, and business meaning at defined boundaries. Do not
equate validity with authorization: syntactically valid invoice ID
\texttt{42} may still belong to another tenant.

\begin{sourceconcept}
Chess and West distinguish security-sensitive data entry points from interpreters and dangerous operations; Dowd et al. emphasize that misplaced trust crosses application and OS interfaces. Safe design removes unintended interpretation rather than trying to guess every hostile input spelling.
\end{sourceconcept}

\subsection{Enforce operation- and object-level
authorization}\label{enforce-operation--and-object-level-authorization}

\begin{itemize}
\tightlist
\item
  Use a server-trusted identity; never use a caller-supplied tenant
  field as the authority by default.
\item
  Enforce permissions at the protected operation and scope data access
  to the principal.
\item
  Check every path: CRUD, search, bulk export, background jobs, cache,
  and GraphQL resolvers.
\item
  Deny by default when authorization context is absent or indeterminate.
\item
  Test cross-role, cross-tenant, disabled-account, and time-of-check
  changes.
\end{itemize}

\section{Python: database and filesystem
boundaries}\label{python-database-and-filesystem-boundaries}

\subsection{Parameter binding}\label{parameter-binding}

\begin{Shaded}
\begin{Highlighting}[]
\CommentTok{\# Vulnerable: user input becomes part of SQL grammar.}
\KeywordTok{def}\NormalTok{ load\_user\_bad(conn, user):}
    \ControlFlowTok{return}\NormalTok{ conn.execute(}\StringTok{"SELECT * FROM users WHERE name=\textquotesingle{}"} \OperatorTok{+}\NormalTok{ user }\OperatorTok{+} \StringTok{"\textquotesingle{}"}\NormalTok{)}

\CommentTok{\# Better: use placeholders supported by your database driver.}
\KeywordTok{def}\NormalTok{ load\_user(conn, user):}
    \ControlFlowTok{return}\NormalTok{ conn.execute(}\StringTok{"SELECT * FROM users WHERE name=?"}\NormalTok{, (user,))}
\end{Highlighting}
\end{Shaded}

DB-API placeholder syntax varies. Parameterization handles
\emph{values}, not arbitrary table names or \texttt{ORDER\ BY}
expressions; construct identifiers from fixed allowlists. Also enforce
least-privilege database credentials.

\subsection{Path confinement}\label{path-confinement}

\begin{Shaded}
\begin{Highlighting}[]
\ImportTok{from}\NormalTok{ pathlib }\ImportTok{import}\NormalTok{ Path}

\KeywordTok{def}\NormalTok{ read\_report(base\_dir: Path, name: }\BuiltInTok{str}\NormalTok{) }\OperatorTok{{-}\textgreater{}} \BuiltInTok{bytes}\NormalTok{:}
\NormalTok{    root }\OperatorTok{=}\NormalTok{ base\_dir.resolve(strict}\OperatorTok{=}\VariableTok{True}\NormalTok{)}
\NormalTok{    candidate }\OperatorTok{=}\NormalTok{ (root }\OperatorTok{/}\NormalTok{ name).resolve(strict}\OperatorTok{=}\VariableTok{True}\NormalTok{)}
    \ControlFlowTok{if} \KeywordTok{not}\NormalTok{ candidate.is\_relative\_to(root):}
        \ControlFlowTok{raise} \PreprocessorTok{PermissionError}\NormalTok{(}\StringTok{"path escapes report root"}\NormalTok{)}
    \ControlFlowTok{return}\NormalTok{ candidate.read\_bytes()}
\end{Highlighting}
\end{Shaded}

This checks the resolved path at one point in time. It is \textbf{not
race-safe} against hostile concurrent symlink replacement. Where an
attacker can modify path components concurrently, use
descriptor-relative/open-at APIs and platform-appropriate confinement,
or keep the directory tree unmodifiable by untrusted principals. Python
\texttt{Path.is\_relative\_to} requires an appropriate supported
interpreter version.

\section{C and C++: memory bounds and integer
range}\label{c-and-c-memory-bounds-and-integer-range}

\subsection{Checked allocation and bounded
copying}\label{checked-allocation-and-bounded-copying}

\begin{Shaded}
\begin{Highlighting}[]
\PreprocessorTok{\#include }\ImportTok{\textless{}stdint.h\textgreater{}}
\PreprocessorTok{\#include }\ImportTok{\textless{}stdlib.h\textgreater{}}
\PreprocessorTok{\#include }\ImportTok{\textless{}string.h\textgreater{}}

\DataTypeTok{int}\NormalTok{ build\_buffer}\OperatorTok{(}\DataTypeTok{const} \DataTypeTok{unsigned} \DataTypeTok{char} \OperatorTok{*}\NormalTok{src}\OperatorTok{,} \DataTypeTok{size\_t}\NormalTok{ count}\OperatorTok{,}
                 \DataTypeTok{size\_t}\NormalTok{ item\_size}\OperatorTok{,} \DataTypeTok{unsigned} \DataTypeTok{char} \OperatorTok{**}\NormalTok{out}\OperatorTok{)} \OperatorTok{\{}
    \ControlFlowTok{if} \OperatorTok{(!}\NormalTok{out }\OperatorTok{||} \OperatorTok{!}\NormalTok{src }\OperatorTok{||} \OperatorTok{!}\NormalTok{item\_size}\OperatorTok{)} \ControlFlowTok{return} \OperatorTok{{-}}\DecValTok{1}\OperatorTok{;}
    \OperatorTok{*}\NormalTok{out }\OperatorTok{=}\NormalTok{ NULL}\OperatorTok{;}
    \ControlFlowTok{if} \OperatorTok{(}\NormalTok{count }\OperatorTok{\textgreater{}}\NormalTok{ SIZE\_MAX }\OperatorTok{/}\NormalTok{ item\_size}\OperatorTok{)} \ControlFlowTok{return} \OperatorTok{{-}}\DecValTok{1}\OperatorTok{;}
    \DataTypeTok{size\_t}\NormalTok{ bytes }\OperatorTok{=}\NormalTok{ count }\OperatorTok{*}\NormalTok{ item\_size}\OperatorTok{;}
    \DataTypeTok{unsigned} \DataTypeTok{char} \OperatorTok{*}\NormalTok{p }\OperatorTok{=}\NormalTok{ malloc}\OperatorTok{(}\NormalTok{bytes }\OperatorTok{?}\NormalTok{ bytes }\OperatorTok{:} \DecValTok{1}\OperatorTok{);}
    \ControlFlowTok{if} \OperatorTok{(!}\NormalTok{p}\OperatorTok{)} \ControlFlowTok{return} \OperatorTok{{-}}\DecValTok{1}\OperatorTok{;}
    \CommentTok{/* Requires caller guarantee: src points to \textgreater{}= bytes readable bytes. */}
\NormalTok{    memcpy}\OperatorTok{(}\NormalTok{p}\OperatorTok{,}\NormalTok{ src}\OperatorTok{,}\NormalTok{ bytes}\OperatorTok{);}
    \OperatorTok{*}\NormalTok{out }\OperatorTok{=}\NormalTok{ p}\OperatorTok{;}
    \ControlFlowTok{return} \DecValTok{0}\OperatorTok{;}
\OperatorTok{\}}
\end{Highlighting}
\end{Shaded}

The size guard prevents multiplication wraparound. It does not establish
that \texttt{src} has \texttt{bytes} readable bytes, that
source/destination overlap is legal, or that ownership is correct. A
robust API must supply an explicit source length and enforce
\texttt{bytes\ \textless{}=\ src\_len}; use memory-safe abstractions
where feasible. Compile with warnings, analyzers, and sanitizers in test
builds.

\subsection{Native checklist}\label{native-checklist}

\begin{enumerate}
\def\labelenumi{\arabic{enumi}.}
\tightlist
\item
  Prove arithmetic before allocation; distinguish signed overflow and
  unsigned wraparound.
\item
  Compare allocation size, source size, and copy length in the same
  units.
\item
  Account for NUL terminators, encoding, and multibyte lengths where
  relevant.
\item
  Verify lifetime, ownership, aliasing, failure cleanup, and integer
  conversion.
\item
  Treat races, signal handlers, file descriptors, and privilege drops as
  part of correctness.
\end{enumerate}

\section{Java / JDBC: interpreter
separation}\label{java-jdbc-interpreter-separation}

\begin{Shaded}
\begin{Highlighting}[]
\BuiltInTok{String}\NormalTok{ sql }\OperatorTok{=} \StringTok{"SELECT id, owner\_id FROM invoices WHERE owner\_id = ?"}\OperatorTok{;}
\ControlFlowTok{try} \OperatorTok{(}\BuiltInTok{PreparedStatement}\NormalTok{ ps }\OperatorTok{=}\NormalTok{ connection}\OperatorTok{.}\FunctionTok{prepareStatement}\OperatorTok{(}\NormalTok{sql}\OperatorTok{))} \OperatorTok{\{}
\NormalTok{    ps}\OperatorTok{.}\FunctionTok{setLong}\OperatorTok{(}\DecValTok{1}\OperatorTok{,}\NormalTok{ authenticatedOwnerId}\OperatorTok{);}
    \ControlFlowTok{try} \OperatorTok{(}\BuiltInTok{ResultSet}\NormalTok{ rows }\OperatorTok{=}\NormalTok{ ps}\OperatorTok{.}\FunctionTok{executeQuery}\OperatorTok{())} \OperatorTok{\{}
        \ControlFlowTok{while} \OperatorTok{(}\NormalTok{rows}\OperatorTok{.}\FunctionTok{next}\OperatorTok{())} \OperatorTok{\{} \CommentTok{/* process authorized records */} \OperatorTok{\}}
    \OperatorTok{\}}
\OperatorTok{\}}
\end{Highlighting}
\end{Shaded}

Prepared statements protect against SQL-grammar injection at bound-value
positions; they do not decide \emph{whose} invoices should be returned.
The \texttt{authenticatedOwnerId} must originate from a validated
server-side principal. Test against wrong-tenant access and missing
privilege, not just SQL metacharacters.

\section{JavaScript / TypeScript: browser and API
boundaries}\label{javascript-typescript-browser-and-api-boundaries}

\begin{Shaded}
\begin{Highlighting}[]
\CommentTok{// Risky: interprets untrusted data as markup.}
\NormalTok{messageElement}\OperatorTok{.}\AttributeTok{innerHTML} \OperatorTok{=}\NormalTok{ response}\OperatorTok{.}\AttributeTok{message}\OperatorTok{;}

\CommentTok{// Safer for plain text:}
\NormalTok{messageElement}\OperatorTok{.}\AttributeTok{textContent} \OperatorTok{=}\NormalTok{ response}\OperatorTok{.}\AttributeTok{message}\OperatorTok{;}
\end{Highlighting}
\end{Shaded}

For rich HTML, use a maintained, context-appropriate sanitization
library and a restrictive content security policy as defense in depth.
Avoid string concatenation into SQL/NoSQL/commands. For server-side
APIs, validate parsed request objects, enforce authorization on every
route, and avoid assuming TypeScript types remain present at runtime.

\section{Go: secure process
invocation}\label{go-secure-process-invocation}

\begin{Shaded}
\begin{Highlighting}[]
\CommentTok{// Avoid shell composition when calling programs.}
\NormalTok{cmd }\OperatorTok{:=}\NormalTok{ exec}\OperatorTok{.}\NormalTok{CommandContext}\OperatorTok{(}\NormalTok{ctx}\OperatorTok{,} \StringTok{"convert"}\OperatorTok{,} \StringTok{"{-}{-}"}\OperatorTok{,}\NormalTok{ inputPath}\OperatorTok{,}\NormalTok{ outputPath}\OperatorTok{)}
\ControlFlowTok{if}\NormalTok{ err }\OperatorTok{:=}\NormalTok{ cmd}\OperatorTok{.}\NormalTok{Run}\OperatorTok{();}\NormalTok{ err }\OperatorTok{!=} \OtherTok{nil} \OperatorTok{\{}
    \ControlFlowTok{return}\NormalTok{ err}
\OperatorTok{\}}
\end{Highlighting}
\end{Shaded}

The example illustrates passing arguments separately. A literal
\texttt{-\/-} works only if the target executable supports that
convention; validate allowed executable options, constrain paths,
timeouts, OS privileges, and output size. A no-shell API reduces shell
metacharacter risk but does not make arbitrary command-line arguments
safe.

\section{Errors, resources, and
secrets}\label{errors-resources-and-secrets}

\textbf{Error handling:} never suppress a security-critical exception
and continue in a privileged state. Return generic externally visible
errors while recording sufficient protected internal context for
troubleshooting. Enforce cleanup of files, connections, locks,
transactions, and temporary credentials on every exit path.

\textbf{Credential lifecycle:} prefer workload identity where possible;
otherwise use approved secret storage with narrow permissions. On
confirmed exposure, revoke/rotate credentials, update dependents, assess
logs and downstream compromise, and eliminate the exposure path. Do not
paste tokens into report screenshots, test fixtures, or CI logs.

\begin{modernextension}
GitHub Secret Protection provides secret scanning and push protection subject to supported patterns, repository entitlement, and configuration. Supported patterns vary by provider and token type; a bypass should trigger a review process. Security controls should be verified against GitHub's current docs and tested with non-sensitive synthetic samples, never live production credentials.
\end{modernextension}

\section{Secure development
integration}\label{secure-development-integration}

\begin{longtable}[]{@{}
  >{\raggedright\arraybackslash}p{(\columnwidth - 4\tabcolsep) * \real{0.3333}}
  >{\raggedright\arraybackslash}p{(\columnwidth - 4\tabcolsep) * \real{0.3333}}
  >{\raggedright\arraybackslash}p{(\columnwidth - 4\tabcolsep) * \real{0.3333}}@{}}
\toprule\noalign{}
\begin{minipage}[b]{\linewidth}\raggedright
Stage
\end{minipage} & \begin{minipage}[b]{\linewidth}\raggedright
Required evidence
\end{minipage} & \begin{minipage}[b]{\linewidth}\raggedright
Practical gate
\end{minipage} \\
\midrule\noalign{}
\endhead
\bottomrule\noalign{}
\endlastfoot
Requirements & Security invariants; sensitive data map & Security story
acceptance criteria \\
Design & Trust boundaries; threat model & Owner approves high-risk
assumptions \\
Implementation & Safe APIs; code review; unit tests & Negative tests in
feature branch \\
Build & Dependency lock, provenance, scanner logs & Successful build and
expected checks \\
PR & CodeQL results where supported & Required branch-protection
checks \\
Release & Config review, secrets posture, rollback & Approval tied to
immutable artifact \\
Operate & Logging, response plan, patch ownership & SLOs and
vulnerability follow-up \\
\end{longtable}

\textbf{Preventive principle:} changes to analysis configuration require
review just like application code. An empty report caused by a broken
extractor or skipped CI job is not success. Record what languages/files
were analyzed and whether the scan succeeded.

\section{Cross-language review
checklist}\label{cross-language-review-checklist}

\begin{itemize}
\tightlist
\item
  Identify security-sensitive entry points and untrusted inputs,
  including background triggers.
\item
  Enforce identity, tenant, permission, and state transitions at
  authoritative boundaries.
\item
  Use parameterized APIs for SQL, no-shell process invocation,
  context-specific output encoding.
\item
  Validate paths with confinement appropriate to attacker concurrency
  capabilities.
\item
  Limit requests, allocations, loops, files, recursion, and response
  sizes.
\item
  Use approved cryptographic libraries and a defined key management
  design; avoid invented cryptography.
\item
  Protect secrets and personal data in logs, exceptions, traces, caches,
  and builds.
\item
  Write regression tests for the \emph{rejected} behavior as well as the
  valid case.
\end{itemize}

\begin{fieldwork}
**Remediation review:** Choose one unsafe data-flow example and write (1) a security invariant, (2) a patch that preserves intended functionality, (3) a negative regression test, (4) a discussion of residual risk, and (5) evidence that the PR scan and build actually covered the affected files.
\end{fieldwork}

\section{Annex A: additional language and framework
patterns}\label{annex-a-additional-language-and-framework-patterns}

\subsection{C\# / ASP.NET Core: policy
enforcement}\label{c-asp.net-core-policy-enforcement}

ASP.NET Core authentication middleware and authorization attributes
provide a useful base, but a role-based controller attribute does not
automatically implement row-level ownership. Treat business
authorization as a service or resource decision after obtaining a
server-authenticated principal.

\begin{Shaded}
\begin{Highlighting}[]
\CommentTok{// Illustrative application{-}layer authorization.}
\KeywordTok{public}\NormalTok{ async Task}\OperatorTok{\textless{}}\NormalTok{Invoice}\OperatorTok{?\textgreater{}} \FunctionTok{GetInvoiceAsync}\OperatorTok{(}
\NormalTok{    ClaimsPrincipal principal}\OperatorTok{,} \DataTypeTok{long}\NormalTok{ invoiceId}\OperatorTok{)} \OperatorTok{\{}
    \DataTypeTok{var}\NormalTok{ tenant }\OperatorTok{=}\NormalTok{ principal}\OperatorTok{.}\FunctionTok{FindFirst}\OperatorTok{(}\StringTok{"tenant"}\OperatorTok{)?.}\FunctionTok{Value}\OperatorTok{;}
    \KeywordTok{if} \OperatorTok{(}\DataTypeTok{string}\OperatorTok{.}\FunctionTok{IsNullOrWhiteSpace}\OperatorTok{(}\NormalTok{tenant}\OperatorTok{))} \KeywordTok{return} \KeywordTok{null}\OperatorTok{;}
    \CommentTok{// Canonical data{-}access query enforces tenant scoping.}
    \KeywordTok{return}\NormalTok{ await db}\OperatorTok{.}\FunctionTok{Invoices}\OperatorTok{.}\FunctionTok{SingleOrDefaultAsync}\OperatorTok{(}\NormalTok{x }\OperatorTok{=\textgreater{}}
\NormalTok{        x}\OperatorTok{.}\FunctionTok{Id} \OperatorTok{==}\NormalTok{ invoiceId }\OperatorTok{\&\&}\NormalTok{ x}\OperatorTok{.}\FunctionTok{TenantId} \OperatorTok{==}\NormalTok{ tenant}\OperatorTok{);}
\OperatorTok{\}}
\end{Highlighting}
\end{Shaded}

The trustworthiness of a \texttt{tenant} claim depends on token
validation, issuer, audience and claim issuance controls; do not blindly
accept a user-edited field. Where permissions change quickly, token
contents alone may be stale. Validate role/action and business state in
addition to tenant scope, as appropriate.

\subsection{Rust: memory-safe does not mean
logic-safe}\label{rust-memory-safe-does-not-mean-logic-safe}

Rust's ownership and borrowing rules prevent many traditional memory
lifetime defects in safe code, but do not prove authorization
correctness, prevent SQL injection from unchecked string construction,
or guarantee safe behavior of \texttt{unsafe} blocks or foreign-language
calls. Use parameterized database APIs, strong types, checked
arithmetic, error propagation, and tests for the expected authorization
invariant. For an \texttt{unsafe} module, document invariants at its
boundary and make callers unable to violate them by accident.

\subsection{Node.js / API input
processing}\label{node.js-api-input-processing}

An HTTP request parser produces dynamically shaped data, even when
TypeScript types are declared for application code. Runtime validation
must reject unexpected types and excessive sizes before expensive
operations. Object property copying into domain models should use
explicit allowlists. For URLs, require a narrowly specified protocol and
destination policy, review DNS and redirects, and impose network egress
restrictions rather than assuming string prefix checks prove safety.

\subsection{Cryptography and tokens}\label{cryptography-and-tokens}

Use maintained cryptographic protocols and libraries. Application
developers should not implement a block cipher or invent a signature
format as a substitute for a trusted standard. For passwords, use a
current, configured password-hashing mechanism appropriate to the
deployment and evolving organizational policy. For tokens, validate
issuer, audience, expiry, algorithm constraints, signature or binding,
and the intended action; protect keys with lifecycle controls and audit
changes.

\section{Annex B: code review as a structured
gate}\label{annex-b-code-review-as-a-structured-gate}

A secure code review should answer four independent questions:
\textbf{Does the code implement the business requirement? Does it
preserve the security invariant? Does it handle failure and concurrency
safely? Did the verification actually run?} A passing scanner can
contribute only to the last two, and then only within its configured
model and scope.

For each high-risk change, ask the author to supply: the abuse case;
vulnerable and secure behavior; a targeted negative test; list of
affected handlers/queues/jobs; data and privilege impact; dependency or
configuration changes; and evidence that the relevant language was
analyzed on the PR SHA.

\textbf{Example review comment:} ``The check verifies a role but not
ownership of \texttt{invoice\_id}. Please add a tenant-scoped query at
the repository boundary and tests for a valid invoice belonging to a
different tenant. Include the export worker that calls the same
repository method.'' Such comments are measurable and reduce repeated
disputes over vague security expectations.

\section{Annex C: defensive pattern
catalog}\label{annex-c-defensive-pattern-catalog}

\begin{longtable}[]{@{}
  >{\raggedright\arraybackslash}p{(\columnwidth - 4\tabcolsep) * \real{0.3333}}
  >{\raggedright\arraybackslash}p{(\columnwidth - 4\tabcolsep) * \real{0.3333}}
  >{\raggedright\arraybackslash}p{(\columnwidth - 4\tabcolsep) * \real{0.3333}}@{}}
\toprule\noalign{}
\begin{minipage}[b]{\linewidth}\raggedright
Failure mechanism
\end{minipage} & \begin{minipage}[b]{\linewidth}\raggedright
Preferred design direction
\end{minipage} & \begin{minipage}[b]{\linewidth}\raggedright
Regression evidence
\end{minipage} \\
\midrule\noalign{}
\endhead
\bottomrule\noalign{}
\endlastfoot
User-controlled DB grammar & Prepared statements and allowlisted
identifiers & Metacharacters remain bound values \\
Shell command composition & No-shell argument API and command allowlist
& Argument cannot become shell syntax \\
User-controlled HTML & Contextual encoding / text-only DOM sink &
Markup-like text remains inert \\
Cross-tenant selection & Authoritative resource-scope query & Valid
foreign resource ID denied \\
Integer width mismatch & Checked range and conversion & Boundary and
over-limit tests \\
Mutable path race & Descriptor-relative atomic operation & Controlled
concurrent replacement test \\
Exposed credential & Revocation, least privilege, approved store & Old
token rejected and scanning verified \\
Duplicate message action & Idempotency and atomic persistent constraint
& Retried message has no second effect \\
Unbounded work & Quota, timeouts, size checks, scheduling & Oversized
request has bounded resource use \\
\end{longtable}

These are \emph{design directions}, not drop-in guarantees; each must be
adapted to the language/runtime, attacker model, and data contract.
Check for failures caused by newly introduced controls such as overly
broad exception swallowing, incorrect caching of authorization
decisions, or unsafe logging of rejected input.

\section{Source and modern extension
map}\label{source-and-modern-extension-map}

\textbf{Book-grounded:} Dowd et al., Chapters 5--8 (memory, integer and
metacharacter handling), Chapters 9--13 (platform and synchronization),
Chapters 17--18 (application risks); Chess \& West, Chapters 5--12
(validation, memory, exceptions, Web, privacy, privileges); Bishop,
Chapters 14--18.

\textbf{Modern extension:} \url{https://owasp.org/projects/asvs};
\url{https://cheatsheetseries.owasp.org/}; \url{https://cwe.mitre.org/};
\url{https://docs.github.com/en/code-security/concepts/secret-security/about-alerts};
\url{https://csrc.nist.gov/pubs/sp/800/218/final}.

\end{document}
